
\documentclass[12pt,a4paper]{report}
\usepackage[utf8]{inputenc}
\usepackage[T1]{fontenc}
\usepackage[brazil]{babel}
\usepackage{geometry}
\usepackage{graphicx}
\usepackage{fancyhdr}
\usepackage{booktabs}
\usepackage{xcolor}
\usepackage{tikz}
\usepackage{pgfplots}
\usepackage{hyperref}
\usepackage{tcolorbox}
\usepackage{enumitem}

\geometry{margin=2.5cm}
\pgfplotsset{compat=1.17}

\pagestyle{fancy}
\fancyhf{}
\fancyhead[L]{\textbf{AEONCOSMA - Relatório Completo}}
\fancyhead[R]{\today}
\fancyfoot[C]{\thepage}

\definecolor{aeonblue}{RGB}{0,102,204}
\definecolor{critical}{RGB}{220,20,60}
\definecolor{warning}{RGB}{255,165,0}
\definecolor{success}{RGB}{34,139,34}

\title{
    \vspace{-2cm}
    \begin{center}
        \textbf{\Huge AEONCOSMA}\\[0.8cm]
        \textbf{\LARGE Relatório Abrangente do Sistema}\\[0.5cm]
        \textcolor{aeonblue}{\Large Digital Twin Network Analysis}\\[0.3cm]
        \textcolor{gray}{\large Período: 26/07/2025 - 02/08/2025}
    \end{center}
    \vspace{-1cm}
}

\date{02 de August de 2025}

\begin{document}

\maketitle
\thispagestyle{empty}

\newpage

\tableofcontents

\newpage

\chapter{Executive Summary}


Este relatório abrangente analisa o desempenho, segurança e operações da rede AEONCOSMA 
durante o período de 26/07/2025 - 02/08/2025. O sistema processou dados de N/A nós 
e detectou 0 padrões de comportamento distintos.

\textbf{Principais Indicadores:}
\begin{itemize}
\item Nível de Segurança: \textcolor{aeonblue}{85.0\% - MUITO BOM}
\item Disponibilidade da Rede: 95.0\%
\item Performance Média: 50.0\% de eficiência
\item Consenso: 90.0\% de participação
\end{itemize}


\chapter{Análise de Segurança}


A análise de segurança revela um ambiente bem protegido com sistemas de detecção ativos. 
O nível geral de segurança é 85.0\%, 
com 0 ameaças detectadas no período.

\section{Métricas de Segurança}

\subsection{Detecções de Ameaças}
O sistema identificou padrões suspeitos em 0 ocasiões, 
todas adequadamente catalogadas e analisadas.

\subsection{Integridade da Rede}
A integridade dos nós permanece alta, com 95.0\% 
de disponibilidade média e 0 falhas de consenso registradas.


\chapter{Performance e Operações}


A análise de performance indica operação eficiente com métricas dentro dos parâmetros esperados.

\section{Métricas Operacionais}

\subsection{Latência de Rede}
\begin{itemize}
\item Latência média: 40.0ms
\item Latência máxima: 120.0ms
\item SLA de 100ms: Não atendido
\end{itemize}

\subsection{Utilização de Recursos}
\begin{itemize}
\item CPU média: 45.0\%
\item Memória média: 60.0\%
\item Eficiência geral: 47.5\%
\end{itemize}


\chapter{Detecção de Anomalias}


O sistema de detecção de anomalias emprega múltiplos algoritmos para identificar comportamentos suspeitos.

\section{Algoritmos de Detecção}

\subsection{Detecção Simbólica}
Analisou 0 padrões simbólicos, 
identificando 3 eventos anômalos.

\subsection{Análise de Entropia}
Score atual: 4.50 
(Normal: 3.0-6.0, Anômalo: <3.0 ou >6.0)

\subsection{Desvios Estatísticos}
2 alertas por desvios das métricas baseline.


\chapter{Stress Tests}

Nenhum teste de stress foi executado no período.

\chapter{Recomendações}


\section{Recomendações Prioritárias}

\subsection{Curto Prazo (1-4 semanas)}
\begin{enumerate}
\item Implementar monitoramento em tempo real de métricas críticas
\item Estabelecer alertas automáticos para desvios significativos
\item Revisar e otimizar configurações de rede
\item Treinar equipe operacional nos novos procedimentos
\end{enumerate}

\subsection{Médio Prazo (1-3 meses)}
\begin{enumerate}
\item Desenvolver sistema de resposta automática a incidentes
\item Implementar redundância geográfica dos nós críticos
\item Estabelecer plano de recuperação de desastres
\item Criar dashboard executivo para tomada de decisões
\end{enumerate}

\subsection{Longo Prazo (3-12 meses)}
\begin{enumerate}
\item Pesquisar e implementar algoritmos de IA para predição de falhas
\item Desenvolver capacidades de auto-healing da rede
\item Estabelecer parcerias para redundância de infraestrutura
\item Implementar sistema de auditoria contínua
\end{enumerate}


\chapter{Apêndices}


\section{Apêndice A: Dados Brutos}


\begin{verbatim}
Timestamp de Coleta: 2025-08-02T16:00:19.602958
Nós Totais: N/A
Nós Online: N/A
CPU Média: N/A%
Memória Média: N/A%
Latência Média: N/Ams
Taxa de Consenso: N/A
\end{verbatim}


\section{Apêndice B: Metodologia}

\subsection{Coleta de Dados}
Os dados foram coletados através de agentes distribuídos em cada nó da rede, 
transmitindo métricas a cada 30 segundos para o sistema central de monitoramento.

\subsection{Análise Estatística}
Utilizamos análise de séries temporais, detecção de outliers e algoritmos de 
machine learning para identificar padrões anômalos.

\subsection{Validação}
Todos os alertas passam por validação cruzada com múltiplos algoritmos antes 
da classificação final.

\section{Apêndice C: Configurações Técnicas}


\begin{verbatim}
Sistema: AEONCOSMA Digital Twin Network
Versão: 1.0.0
Algoritmo de Consenso: Proof of Stake Modificado
Criptografia: SHA-256, ECDSA
Protocolos de Rede: TCP/IP, UDP, WebSocket
Monitoramento: Real-time com detecção de anomalias
\end{verbatim}


\section{Apêndice D: Glossário}

\begin{description}
\item[AEONCOSMA] Advanced Earth Observation Network for Cosmic Monitoring and Analysis
\item[Digital Twin] Representação digital em tempo real do sistema físico
\item[Consenso] Mecanismo de acordo distribuído entre nós da rede
\item[Entropia] Medida de aleatoriedade ou desordem no sistema
\item[Hash] Função criptográfica de verificação de integridade
\end{description}


\end{document}
